\section*{Supplementary Material}

\setcounter{table}{0}
\renewcommand{\thetable}{S\arabic{table}}

\subsection*{Rod Geometry and Kinematic Parameterization}
The director $\mathbf{d}_3(s,t)$ aligns with the tangent to the centerline in the absence of shear, while $\mathbf{d}_1(s,t)$ and $\mathbf{d}_2(s,t)$ define the principal axes of the vertebral cross-section. The kinematics are governed by the evolution equations:
\begin{align}
\mathbf{r}'(s) &= \mathbf{v}(s) = \mathbf{d}_3(s) + \boldsymbol{\varepsilon}(s), \\
\mathbf{d}_i'(s) &= \mathbf{u}(s) \times \mathbf{d}_i(s),
\label{eq:cosserat_kinematics}
\end{align}
where $\mathbf{v}(s)$ is the tangent vector, $\boldsymbol{\varepsilon}(s)$ represents shear and extension strains, and $\mathbf{u}(s) = \kappa_1 \mathbf{d}_1 + \kappa_2 \mathbf{d}_2 + \tau \mathbf{d}_3$ is the Darboux vector encoding the flexural curvature components $\kappa_{1,2}$ and torsional twist $\tau$.

\subsection*{Cosserat force and moment balance}
For a static rod subject to gravity $\mathbf{f}_g = \rho A \mathbf{g}$ and active moments induced by the information field gradients, the balance equations for internal force $\mathbf{n}$ and moment $\mathbf{m}$ are:
\begin{align}
\mathbf{n}'(s) + \mathbf{f}_g &= \mathbf{0}, \nonumber \\
\mathbf{m}'(s) + \mathbf{r}'(s) \times \mathbf{n}(s) + \mathbf{m}_{\mathrm{info}}'(s) &= \mathbf{0},
\label{eq:cosserat_balance}
\end{align}
where $\mathbf{m}_{\mathrm{info}}$ represents the active couple.

To model the human spine, we impose clamped-free boundary conditions. The sacral base ($s=0$) is rigidly fixed, while the cranial tip ($s=L$) is free of external loads:
\begin{align}
\mathbf{r}(0) = \mathbf{0}, \quad &\mathbf{d}_i(0) = \mathbf{e}_i \quad (\text{Clamped Base}) \nonumber \\
\mathbf{n}(L) = \mathbf{0}, \quad &\mathbf{m}(L) = \mathbf{0} \quad (\text{Free Tip})
\label{eq:boundary_conditions}
\end{align}
where $\{\mathbf{e}_i\}$ is the laboratory frame.

\begin{table}[h!]
  \centering
  \caption{\textbf{The Biological Equivalence Principle.} Correspondence between the geometric variables of General Relativity (GR) and the morphoelastic variables of the Information-Energy-Countercurvature (IEC) framework. This analogy motivates the mathematical structure of the cost functional in Eq.~\ref{eq:cost_functional}.}
  \label{tab:gr_analogy}
  \small
  \begin{tabular}{p{3cm} p{5cm} p{5cm}}
    \toprule
    \textbf{Concept} & \textbf{General Relativity (GR)} & \textbf{IEC Framework} \\
    \midrule
    \textit{Geometry} & Metric tensor $g_{\mu\nu}$ & Reference configuration $\hat{\mathbf{u}}(s)$ \\
    \textit{Source} & Stress-Energy tensor $T_{\mu\nu}$ & Mechanical load vector $\mathbf{n}(s)$ \\
    \textit{Equation} & $R_{\mu\nu} - \frac{1}{2}Rg_{\mu\nu} = \kappa T_{\mu\nu}$ & $\mathbf{n}' + \mathbf{f} = \mathbf{0}, \mathbf{m}' + \mathbf{u}' \times \mathbf{n} = \mathbf{0}$ \\
    \textit{Path} & Geodesic ($\delta \int ds = 0$) & Homeostatic shape ($\delta \mathcal{E} = 0$) \\
    \textit{Information} & Information Term $I_{\mu\nu}$ & Information Field $I(s)$ \\
    \textit{Function} & Modulates space curvature & Modulates stress-free strain $\hat{\kappa}$ \\
    \bottomrule
  \end{tabular}
\end{table}
